\section{10 大原创能力深度解析}\label{sec:capabilities-10}

每个能力都是一等公民模块, 独立可测试 / 可调用 / 可解释。每条都有具体技术实现, 不是概念堆叠。

\subsection{C1 Persona Vector — 人设向量化 + Git for Voice}\label{sec:cap-persona}

\textbf{解决的痛点}: 「跟了一次热点后, 粉丝吐槽『你变营销号了』。需要边界:什么热点能跟, 怎么跟才不崩人设。」

\textbf{核心思想}: 人设不是 prompt, 而是一个 \textbf{持续学习的向量}。

\subsubsection{算法栈}

\begin{itemize}
    \item \textbf{特征抽取}: 启发式 (句长 / emoji 密度 / 第一人称频率 / 专业词命中) + LLM 抽取 (formality / technicality / humor\_density 等 10 维)
    \item \textbf{Embedding}: 字符 trigram hash 到 256 维向量 (兜底); 可替换为 sentence-transformers
    \item \textbf{EMA 更新}: $v_{t+1} = (1-\alpha) v_t + \alpha v_{new}$, 默认 $\alpha = 0.15$
    \item \textbf{Drift 检测}: cosine similarity, 阈值 \texttt{warn=0.65, block=0.4}
    \item \textbf{Git for Voice}: 主线 \texttt{main} + 实验分支 \texttt{experimental\_humor}, \texttt{merge}/\texttt{rollback} 操作
\end{itemize}

\subsubsection{API}

\begin{lstlisting}[language=bash]
POST /api/v2/persona/calibrate    # 用 5+ 历史样本初始化
POST /api/v2/persona/ingest       # 增量更新 + Drift 检测
POST /api/v2/persona/branch       # 开实验分支
POST /api/v2/persona/merge        # 合并实验回主线
POST /api/v2/persona/drift_check  # 发布前预检
\end{lstlisting}

\subsection{C2 Replay Graph — 可回放执行图}

\textbf{解决的痛点}: 「KOC 问 \textit{『为什么推荐这个话题』}, 系统能展开完整思考链。」

\textbf{核心思想}: 每一步决策被记录为 DAG 节点, 用 Merkle 链防篡改。

\subsubsection{ReplayNode 数据结构}

\begin{lstlisting}[language=Python]
class ReplayNode:
    node_id: str
    parent_ids: List[str]
    phase: CognitivePhase  # sense/think/plan/act/observe/reflect/update
    actor: str             # 哪个工具/模块产生
    input_hash: str        # 输入 SHA256
    input_summary: str
    output_summary: str
    rejected_alternatives: List[str]  # 否决的备选
    duration_ms: int
    timestamp: datetime
    metadata: Dict
    
    def merkle_hash(self) -> str:
        # 含父节点哈希, 形成链
        return sha256(node_id|actor|phase|input_hash|output_summary|parents_hash)
\end{lstlisting}

\subsubsection{第三方可验证}

\texttt{POST /api/v2/replay/\{run\_id\}} 返回完整 DAG + \texttt{merkle\_valid} 字段。第三方可重新计算 Merkle 链验证未被篡改。

\subsection{C3 SimPredictor — 虚拟受众 A/B}

\textbf{解决的痛点}: 「同一选题 3 版开头, 哪版会火?不能每条都跑真实 A/B test。」

\textbf{核心思想}: 用 LLM 模拟 5 个不同人格 (Generative Agents 启发), 让它们扮演真实读者评分。

\subsubsection{默认 Cohort 库}

\begin{tabularx}{\textwidth}{>{\bfseries}l X}
    \toprule
    young\_female\_t2 & 25-30 岁二线女性 / 美妆护肤 \\
    male\_student & 20 岁男大学生 / 数码科技 \\
    career\_woman\_t1 & 28-35 岁一线职场女性 / 自我成长 \\
    hobby\_uncle & 35-50 岁中产男性 / 兴趣爱好 \\
    mom\_t3 & 30-40 岁妈妈 / 母婴亲子 \\
    \bottomrule
\end{tabularx}

\subsubsection{算法}

\begin{enumerate}
    \item 给每个 cohort 一份人格卡 (年龄 / 兴趣 / 平台习惯)
    \item LLM 扮演该人格, 严格 JSON 输出 \{score, completion\_estimate, controversy, reason\}
    \item Platt scaling 校准分数 (避免 LLM 偏积极)
    \item 多 cohort 聚合 → 加权平均 + 不确定性区间
    \item 多版本赛马 → 输出排序 + 推荐版本
\end{enumerate}

\subsection{C4 Risk-Reward Scorer — 选题套利评分}

\textbf{解决的痛点}: 「这个选题潜力大但风险也高, KOC 自己拿捏不准。」

\textbf{核心思想}: 每个选题输出 \textbf{二维分数} 而不是单维, 让 KOC 像看股票 K 线一样做决策。

\subsubsection{特征工程}

\begin{itemize}
    \item 早期信号强度 (来自 Oracle 的 normalized rank)
    \item 监管敏感词命中 (金融 / 医疗 / 政治 / 未成年保护)
    \item 历史同类事故 (向量检索类似话题的翻车记录)
    \item 对立叙事密度 (Forum Debate 中的反对意见数)
    \item 关键词命中: 数据反差 / 亲身经历 / 智商税 / 翻车 → 提升潜力
\end{itemize}

\subsubsection{输出}

\begin{lstlisting}[language=json]
{
  "potential_score": 0.78,
  "risk_score": 0.41,
  "regulatory_sensitivity": "mid",
  "recommendation": "hedge",
  "reasoning": "数据反差强,但涉及金融,建议做解读不押注",
  "calibrated": true
}
\end{lstlisting}

\subsection{C5 Trend Chain Analyzer — 因果链推理}

\textbf{解决的痛点}: 「KOC 看到的热搜都已经红海, 怎么提前看到下一个?」

\textbf{核心思想}: 建模 \textbf{资本流 → 国际舆论 → 中文媒体 → 大众化梗} 四阶段传导规律, 用 Granger 因果 (弱形式) + LLM 机制叙述。

\subsubsection{典型链条}

\begin{quote-box}
某 AI 公司 \$2B 融资 (Day 0) → HN 前排 / 国际推文热议 (Day 3-7) → 36 氪 / 量子位等中文科技媒体 (Day 7-14) → 抖音/B站玩梗 (Day 14-30)
\end{quote-box}

KOC 在 Day 5 介入做 \textbf{科普向}内容, 比 Day 20 介入做 \textbf{玩梗向} 内容能领先国内 70\% 同行, 且不需要硬蹭。

\subsubsection{API}

\begin{lstlisting}[language=bash]
POST /api/v2/trend/chain
{
  "seed_event": "DeepSeek-V4 发布",
  "platform_signals": {"hackernews": 1, "weibo": 0}
}
\end{lstlisting}

\subsection{C6 Cross-Platform Translator — 跨平台改写}

\textbf{解决的痛点}: 「英文 HN 帖很有料, 但翻译成中文不是『直译』, 是『信息再设计』。」

\textbf{核心思想}: 每个平台有独立的 \textbf{语法模板} (语气 / 长度 / 钩子位置 / 禁忌), 同一信息走不同模板得到不同形态。

\subsubsection{平台语法库}

\begin{tabularx}{\textwidth}{>{\bfseries}l l X}
    \toprule
    平台 & 长度 & 关键模板 \\
    \midrule
    小红书 & 300-600 字 & 封面金句 + 信息图清单 + 评论区引导 \\
    抖音 & 60-100 字脚本 & 15s 钩子 + 反转 + 信息密度 + CTA \\
    视频号 & 90-180s 脚本 & 信任开场 + 慢节奏 + 熟人圈推荐感 \\
    B 站 & 5-15 分钟 & 标题党钩子 + 系统讲解 + 弹幕引导 \\
    微信公众号 & 1500-3000 字 & 新闻钩子 + 多段深度 + 数据论证 \\
    微博 & 140-300 字 & 话题标签 + 短论断 + @ 相关账号 \\
    \bottomrule
\end{tabularx}

\subsection{C7 Cohort Affinity Estimator — 分人群吸引力}

\textbf{解决的痛点}: 「为什么这条内容男大学生爱看, 25 岁女性不感兴趣?」

\textbf{核心思想}: 内容 × 受众卡的二维偏好估计, 复用 SimPredictor 但接口更细。

\textbf{破圈建议}: 当某 cohort 评分明显低于其他时, 推荐「加入 X cohort 共鸣点或换一个开头」。

\subsection{C8 Campaign Planner — 7 天战役图}

\textbf{解决的痛点}: 「KOC 一周 5 更要不透支人设, 需要系列感。」

\textbf{核心思想}: 一周内容规划 = 多约束最优化问题。

\subsubsection{约束}

\begin{itemize}
    \item \textbf{流量结构}: 引流 20\% / 信任 50\% / 转化 30\% (千粉冷启动调为 40/50/10)
    \item \textbf{角色覆盖}: hook / deepdive / qa / ugc / summary / live\_pre / review 必须各至少 1 天
    \item \textbf{热点缓冲}: 至少 1 天 \texttt{hotspot\_buffer} 预留给突发热点
    \item \textbf{依赖}: \texttt{qa} 依赖前 2 天的 hook 或 deepdive
    \item \textbf{库存}: 用户每日可投入小时数 (默认 2 h)
\end{itemize}

\textbf{算法}: LLM 起草 + 贪心调度 + 约束验证。

\subsection{C9 Multi-Modal Coherence — 多模态一致性 Manifest}

\textbf{解决的痛点}: 「文案改了, 视频字幕和 BGM 没跟上, 内容割裂。」

\textbf{核心思想}: 文 / 图 / 视频脚本 / 字幕 / BGM / 封面 共享一个 JSON Schema (\texttt{Manifest}), diff 引擎自动标记需重剪段落。

\subsection{C10 Self-Critique Loop — 自我批评闭环}

\textbf{解决的痛点}: 「LLM 一次写出的文案有时风险高, 需要二次过一遍。」

\textbf{核心思想}: 仿 SELF-REFINE 论文。生成→批评→修订循环, 直到 \textbf{Critic 给出 0 条建议} 或达到 \texttt{max\_rounds=3}。

\subsubsection{停止条件}

\begin{enumerate}
    \item Critic 输出 \texttt{issues=[]}
    \item 达到最大轮数
    \item 能量函数 $E = |issues|/5$ 低于阈值 0.2
\end{enumerate}

\subsubsection{评委交互}

最终展示给评委的是 \textbf{最终稿 + 迭代次数}, 详细 trace 在 Replay 里。

\section{反常识设计 5 条 (评委记忆点)}

\begin{enumerate}
    \item \textbf{故意不直接给唯一答案, 给 3 个对立假设}: KOC 拥有 ownership, 选择即标注, 有利于 Persona 学习
    \item \textbf{明确展示 Agent 的不确定性}: \textit{「我不确定, 但倾向 A, 因为...」}更可信 (心理学:可辩驳性 = 可信)
    \item \textbf{故意对高仿爆款降权}: 鼓励差异 5\%-15\% 的版本, 同质内容边际收益递减
    \item \textbf{默认低幻觉}: 宁少一句狠话, 也要带来源, 短期爽文 < 长期信任
    \item \textbf{每周 60 秒人设校准 ritual}: 避免风格被单次爆款带偏, 同时形成留存仪式
\end{enumerate}

\begin{insight}
这 5 条反常识看起来牺牲了 \textit{「酷炫的一键解决」} 体验, 但对 KOC 的 \textbf{长期增长}是正确的。Ripple 不是替代 KOC 决策, 而是把 KOC 决策变得 \textbf{可解释 / 可复盘 / 可持续}。
\end{insight}
